\begin{table*}[t]
\centering
\caption{Instâncias: número de hexágonos H3 elegíveis, número de pares origem--destino da matriz de impacto (identificador \texttt{impacts} do conjunto de dados), valor da linha de base e melhor solução conhecida (BKS) entre todas as execuções.}
\label{tab:instances}
\footnotesize
\begin{tabular}{@{}llrrrrl@{}}
\toprule
Localidade & Perfil & Hexágonos & Pares origem--destino & $\overline{\IQC}_b$ & BKS & BKS obtida por \\
\midrule
Av. Paulista & atleta & 177 & 4.998 & 54,00 & 63,88 & PSO \\
Av. Paulista & adulto & 177 & 2.999 & 47,06 & 55,05 & PSO \\
Av. Paulista & idoso & 177 & 1.328 & 38,68 & 46,24 & GRASP \\
Itajubá & atleta & 167 & 3.843 & 30,23 & 56,24 & PSO \\
Itajubá & adulto & 167 & 2.273 & 25,41 & 52,73 & PSO \\
Itajubá & idoso & 167 & 1.077 & 22,75 & 46,67 & PSO \\
Ana Rosa & atleta & 177 & 4.904 & 42,54 & 58,22 & PSO \\
Ana Rosa & adulto & 177 & 2.935 & 38,53 & 51,42 & PSO \\
Ana Rosa & idoso & 177 & 1.308 & 32,45 & 38,54 & PSO \\
Milão & atleta & 183 & 7.119 & 62,95 & 66,28 & PSO \\
Milão & adulto & 183 & 4.440 & 57,45 & 60,68 & GRASP \\
Milão & idoso & 183 & 2.123 & 49,55 & 52,88 & GRASP \\
Shinjuku & atleta & 194 & 6.907 & 57,10 & 62,43 & PSO \\
Shinjuku & adulto & 194 & 4.221 & 51,27 & 56,95 & PSO \\
Shinjuku & idoso & 194 & 1.980 & 42,23 & 47,92 & GRASP \\
\bottomrule
\end{tabular}
\end{table*}
